\documentclass[11pt]{article}
\usepackage[T1]{fontenc}
\usepackage{amsmath,amssymb,amsthm}
\usepackage[margin=1in]{geometry}
\usepackage{booktabs}
\usepackage[hidelinks]{hyperref}
\hypersetup{
  pdftitle={The exact generalized-pancake spectral gap at (m,n)=(4,2)},
  pdfauthor={Xiaojian Zeng},
  pdfsubject={Exact spectra of a generalized pancake Cayley graph and its Schreier graph},
  pdfkeywords={generalized pancake graph, spectral gap, Cayley graph, Schreier graph, finite Fourier analysis}
}

\newtheorem{theorem}{Theorem}
\newcommand{\Z}{\mathbb Z}
\newcommand{\C}{\mathbb C}
\newcommand{\Cay}{\mathrm{Cay}}
\newcommand{\Sch}{\mathrm{Sch}}

\title{The exact generalized-pancake spectral gap at $(m,n)=(4,2)$}
\author{Xiaojian Zeng\\
PBC School of Finance, Tsinghua University, 100190\\
\texttt{xiaojianzeng@foxmail.com}}
\date{}

\begin{document}
\maketitle

\begin{abstract}
For the generalized pancake Cayley graph on
$\Z_4\wr S_2$ and its color--position Schreier multigraph, we compute the
unnormalized adjacency spectral gaps exactly.  A Fourier decomposition gives
sixteen explicit $2\times2$ blocks for the $32$-vertex Cayley operator and
four such blocks for the $8$-vertex Schreier operator.  Their characteristic
polynomials show that both second adjacency eigenvalues are $2\sqrt2$.
Consequently
\[
  \psi_{4,2}=\beta_{4,2}=4-2\sqrt2.
\]
This resolves only the single parameter cell $(4,2)$; it makes no assertion
about any other case of the proposed equality.
\end{abstract}

\section{Context and statement}

For $m\geq 3$, a generalized prefix reversal $r_i^\varepsilon$
($\varepsilon\in\{-1,+1\}$) reverses the first $i$ positions of a colored
permutation and adds $\varepsilon$ modulo $m$ to the colors in that prefix.
These reversals generate generalized pancake Cayley graphs; see
\cite{BlancoBuehrle2025,GreavesZhu2026,Blanco2026} for their spectral
context.  Greaves and Zhu conjectured equality between the adjacency gap of
the Cayley graph and that of its color--position quotient
\cite[Conjecture~2]{GreavesZhu2026}.  Blanco subsequently exhibited strict
inequality at $(6,2)$ and proposed a revised parameter classification
\cite[Section~5]{Blanco2026}.  The latter article evaluates $(3,2)$ but does
not state the exact value at $(4,2)$.

We now fix that one cell.  Let
\[
 G=\Z_4\wr S_2,\qquad
 R=\{r_1^-,r_1^+,r_2^-,r_2^+\}.
\]
A vertex of $\Cay(G,R)$ is a two-letter colored permutation, so the Cayley
graph has $2!\,4^2=32$ vertices.  Its generator-sum adjacency matrix is
\[
 A_{\rm Cay}(x,y)=\#\{r\in R:r(x)=y\}.
\]
For the color--position action on $X=\Z_4\times\{1,2\}$, set
\[
 r_i^\varepsilon(a,p)=
 \begin{cases}
   (a+\varepsilon,i+1-p),&p\leq i,\\
   (a,p),&p>i.
 \end{cases}
\]
Its Schreier generator-sum matrix is
\[
 A_{\rm Sch}(u,v)=\#\{r\in R:r(u)=v\}.
\]
Thus fixed generators are counted as loops.  Both matrices are real
symmetric and have row sum four.  If the eigenvalues of a symmetric matrix
$A$ are ordered with multiplicity as
$\lambda_1(A)\geq\lambda_2(A)\geq\cdots$, define
\[
 \psi_{4,2}=4-\lambda_2(A_{\rm Cay}),\qquad
 \beta_{4,2}=4-\lambda_2(A_{\rm Sch}).
\]

\begin{theorem}\label{thm:main}
For the operators just defined,
\[
 \lambda_2(A_{\rm Cay})=\lambda_2(A_{\rm Sch})=2\sqrt2.
\]
Hence $\psi_{4,2}=\beta_{4,2}=4-2\sqrt2$.
\end{theorem}

\section{Exact Fourier blocks}

Put $\zeta=i$ and, with subscripts taken modulo four,
\[
 d_j=\zeta^j+\zeta^{-j}=2\cos(\pi j/2),\qquad
 (d_0,d_1,d_2,d_3)=(2,0,-2,0).
\]
We use only the elementary Fourier basis of the cyclic group $\Z_4$; general
background on finite-group Fourier decompositions may be found in
\cite{Terras1999}.

\subsection{The Schreier operator}

The defining action gives, for $f:X\to\C$,
\begin{align*}
 (A_{\rm Sch}f)(a,1)
   &=\sum_{\varepsilon=\pm1}
        \bigl(f(a+\varepsilon,1)+f(a+\varepsilon,2)\bigr),\\
 (A_{\rm Sch}f)(a,2)
   &=2f(a,2)+\sum_{\varepsilon=\pm1}f(a+\varepsilon,1).
\end{align*}
On functions of the form $f(a,p)=\zeta^{ja}u_p$, the matrix in the position
basis $(1,2)$ is therefore
\[
 S_j=\begin{pmatrix}d_j&d_j\\d_j&2\end{pmatrix}.
\]
The four Fourier modes are mutually orthogonal and exhaust $\C^X$.  Direct
determinant calculation gives, for $j=0,1,2,3$ respectively,
\[
 \det(xI-S_j)=x(x-4),\quad x(x-2),\quad x^2-8,\quad x(x-2).
\]
Consequently
\begin{equation}\label{eq:sch-char}
 \chi_{\rm Sch}(x)=x^3(x-4)(x-2)^2(x^2-8).
\end{equation}

\subsection{The Cayley operator}

Represent a colored word by $(t,a,b)$, where $t=0$ denotes label order
$(1,2)$, $t=1$ denotes label order $(2,1)$, and $a,b\in\Z_4$ are the colors
in the first and second positions.  The four transitions are
\[
 r_1^\varepsilon(t,a,b)=(t,a+\varepsilon,b),\qquad
 r_2^\varepsilon(t,a,b)=(1-t,b+\varepsilon,a+\varepsilon).
\]
Because $R=R^{-1}$, the generator-sum action on functions may equivalently
be written using these forward transitions.  Define
\[
 F_{s,k,l}(t,a,b)=\mathbf 1_{t=s}\zeta^{ka+lb}.
\]
Substitution into the four transitions yields
\begin{equation}\label{eq:intertwine}
 A_{\rm Cay}F_{s,k,l}
   =d_kF_{s,k,l}+d_{k+l}F_{1-s,l,k}.
\end{equation}
It follows that, for $(k,l)\in\Z_4^2$,
\[
 W_{k,l}=\operatorname{span}\{F_{0,k,l},F_{1,l,k}\}
\]
is invariant and has the exact block
\[
 B_{k,l}=\begin{pmatrix}d_k&d_{k+l}\\d_{k+l}&d_l\end{pmatrix}.
\]
The first basis vectors of the sixteen spaces give every Fourier character
on the $t=0$ slice exactly once; the second basis vectors do the same on the
$t=1$ slice.  Thus the spaces are mutually orthogonal and account for all
$32$ dimensions.

For completeness, the determinant factors of all blocks are displayed
below; rows are indexed by $k$ and columns by $l$:
\[
\begin{array}{c|cccc}
 &0&1&2&3\\ \hline
0&x(x-4)&x(x-2)&x^2-8&x(x-2)\\
1&x(x-2)&(x-2)(x+2)&x(x+2)&(x-2)(x+2)\\
2&x^2-8&x(x+2)&x(x+4)&x(x+2)\\
3&x(x-2)&(x-2)(x+2)&x(x+2)&(x-2)(x+2)
\end{array}
\]
Multiplying the sixteen factors gives
\begin{equation}\label{eq:cay-char}
 \chi_{\rm Cay}(x)
 =x^{10}(x-4)(x-2)^8(x+2)^8(x+4)(x^2-8)^2.
\end{equation}

\subsection{Ordering the eigenvalues}

Equations \eqref{eq:sch-char} and \eqref{eq:cay-char} show that the largest
three distinct eigenvalues of each operator are
\[
 4>2\sqrt2>2.
\]
These inequalities are exact: all terms are positive and their squares are
$16>8>4$.  In particular, the factor $x-4$ occurs once in each
characteristic polynomial, and the next eigenvalue, counted with
multiplicity, is $2\sqrt2$.  Subtracting it from the common degree four
proves Theorem~\ref{thm:main}.

\section{Verification and scope}

The proof above is self-contained; computation is not needed for its logical
conclusion.  Two deterministic Python programs were nevertheless used as
cross-checks.  One constructs the $32\times32$ and $8\times8$ integer
matrices directly from the transitions and factors their characteristic
polynomials.  The other verifies equation \eqref{eq:intertwine} and its
Schreier analogue entry by entry over the Gaussian integers, checks block
exhaustion, and compares the block products with direct characteristic
polynomials.  These checks use SymPy \cite{MeurerEtAl2017} and no
floating-point eigenvalues.  Related computer-assisted work from the
user-supplied bibliography likewise records exact rational or algebraic sign
decisions together with independent definition-level verification programs
\cite{ZengEtAl2026}; here, however, the programs are only corroborative.

The result is confined to $(m,n)=(4,2)$.  In particular, it neither proves
the revised classification proposed in \cite{Blanco2026} nor gives an exact
formula for other generalized pancake graphs.  No claim of bibliographic
priority is made.

\section*{Computational and AI assistance disclosure}

OpenAI Codex was used in proof reconstruction, exact-code development,
literature triage, and organization and drafting of the manuscript.  The
author reviewed the mathematical statements, source comparisons, exact
certificates, and final text and takes responsibility for them.  AI systems
were not treated as authors or as sources of mathematical authority.  The
deterministic checks described above are corroboration of the displayed
proof, not a substitute for it.

\bibliographystyle{plain}
\bibliography{references}

\end{document}
